\chapter{نتیجه‌گیری و کارهای آینده}

این فصل جمع‌بندی نهایی پروژه و مسیر توسعه بعدی را ارائه می‌کند. هدف این بخش، تکرار فصل‌های قبل نیست؛ بلکه می‌خواهد روشن کند پروژه امروز کجا ایستاده و قدم‌های بعدی آن چیست.

\section{جمع‌بندی کار انجام‌شده}
\lr{IranAPI Hub} یک پلتفرم اشتراک‌گذاری \lr{API} با بک‌اند \lr{Django REST Framework} و فرانت‌اند \lr{React} است. سامانه کاتالوگ، مستندات، احراز هویت، \lr{Dashboard}، مصرف، اشتراک، سازمان، انتشار، \lr{Caller} و \lr{Studio} را پوشش می‌دهد.

در طول گزارش، پروژه از زاویه کاربر و معماری بررسی شد. فصل نیازمندی‌ها نقش‌ها و جریان‌ها را مشخص کرد. فصل معماری نشان داد این جریان‌ها چگونه بین فرانت‌اند، بک‌اند و داده تقسیم می‌شوند. فصل پیاده‌سازی رفتار بخش‌های اصلی را توضیح داد و فصل آزمون نشان داد بک‌اند با \lr{49} آزمون موفق پشتیبانی می‌شود.

\section{دستاوردهای اصلی}
دستاوردهای اصلی پروژه شامل پشتیبانی از زبان فارسی، مسیرهای عمومی و خصوصی منظم، خروجی \lr{OpenAPI}، تست‌های بک‌اند، ساختار \lr{Docker Compose} و تفکیک نسبی بین لایه نمای \lr{API}، سریال‌سازی، داده و رابط کاربری است.

دستاورد مهم دیگر، هماهنگی نسبی بین تجربه کاربر و معماری فنی است. کاربر کاتالوگ، مستندات، حساب و ابزارهای توسعه‌دهنده را می‌بیند و پشت آن، \lr{API}، \lr{Serializer}، لایه داده و آزمون قرار گرفته‌اند. این رابطه باعث می‌شود پروژه فقط نمایشی نباشد و از نظر فنی هم قابل ادامه‌دادن باشد.

\section{محدودیت‌های فعلی}
محدودیت‌های مستندشده عبارت‌اند از: آماده نبودن تنظیمات تولید، باقی‌ماندن مسیرها و مدل‌های سازگار قدیمی، نبود اجرای تازه آزمون‌های فرانت‌اند در این گزارش، نبود اثبات اتصال واقعی \lr{Caller} به سرویس خارجی و نیاز ورود اجتماعی به پیکربندی ارائه‌دهنده.

این محدودیت‌ها طبیعی‌اند، اما باید در برنامه ادامه کار جدی گرفته شوند. تا زمانی که تنظیمات تولید، رازها، پایگاه داده واقعی، پرداخت، اتصال بیرونی و آزمون مرورگر کامل نشده باشند، بهتر است پروژه به‌عنوان نسخه آموزشی یا نمونه قابل توسعه معرفی شود، نه محصول تولیدی آماده بهره‌برداری عمومی.

\section{کارهای آینده}
کارهای آینده پیشنهادی:
\begin{itemize}
  \item انتقال رازها و تنظیمات حساس به متغیرهای محیطی و غیرفعال‌کردن \lr{DEBUG} در تولید.
  \item تصمیم‌گیری نهایی درباره مرز \lr{MongoDB} و \lr{Django ORM} و حذف دوگانگی غیرضروری.
  \item اجرای منظم آزمون‌های \lr{Playwright} و افزودن آن‌ها به فرایند \lr{CI}.
  \item اتصال \lr{Caller} به سرویس‌های مقصد با کنترل امنیت، \lr{timeout} و ثبت خطای واقعی.
  \item تکمیل پرداخت واقعی، \lr{webhook}ها و کنترل سهمیه مصرف.
  \item افزودن مستندات نصب تولید، نسخه‌گذاری \lr{API} و مانیتورینگ عملیاتی.
\end{itemize}

ترتیب انجام این کارها بهتر است از ریسک‌های بنیادی شروع شود. ابتدا باید امنیت پیکربندی و رازها اصلاح شود، سپس تصمیم داده‌ای بین \lr{MongoDB} واقعی و مدل‌های \lr{ORM} نهایی شود. پس از آن، آزمون‌های مرورگر، پرداخت واقعی، مانیتورینگ و اتصال \lr{Caller} می‌توانند با اطمینان بیشتری کامل شوند.

\section{جمع‌بندی فصل}
پروژه پایه فنی مناسبی برای یک بازارچه \lr{API} دارد و آزمون‌های بک‌اند موجود آن با موفقیت اجرا شده‌اند. مسیر توسعه بعدی باید روی سخت‌سازی تولیدی، یکپارچه‌سازی داده و آزمون‌های انتهابه‌انتها متمرکز شود. با انجام این گام‌ها، سامانه می‌تواند از نمونه آموزشی کامل به محصولی پایدارتر و آماده‌تر برای استفاده واقعی تبدیل شود.
